\section*{अध्याय \ref{ch:b1:complex} --- सम्मिश्र संख्याएँ}

\begin{solution}{exo:b1:complex:1}
$1 + \iu = \sqrt 2\, \eu^{\iu\pi/4}$;
$\;\sqrt 3 - \iu = 2\, \eu^{-\iu\pi/6}$;
$\;-5 = 5\, \eu^{\iu\pi}$;
\[
\frac{1 + \iu}{\sqrt 3 - \iu}
= \frac{\sqrt 2}{2}\, \eu^{\iu(\pi/4 + \pi/6)}
= \frac{\sqrt 2}{2}\, \eu^{5\iu\pi/12}.
\]
वही भागफल बीजगणितीय रूप से संगणित करना (\omterm{def:b1:complex:field}{संयुग्मी} से गुणा कीजिए):
\[
\frac{(1 + \iu)(\sqrt 3 + \iu)}{4}
= \frac{(\sqrt 3 - 1) + \iu(\sqrt 3 + 1)}{4}.
\]
$\frac{\sqrt 2}{2}(\cos\frac{5\pi}{12} + \iu\sin\frac{5\pi}{12})$ से पहचान
करने पर:
\[
\cos\frac{5\pi}{12} = \frac{\sqrt 6 - \sqrt 2}{4},
\qquad
\sin\frac{5\pi}{12} = \frac{\sqrt 6 + \sqrt 2}{4}.
\]
\end{solution}

\begin{solution}{exo:b1:complex:2}
$(1+\iu)^2 = 2\iu$, अतः $(1+\iu)^{20} = (2\iu)^{10} = 2^{10}\,\iu^{10} =
1024 \times (-1) = -1024$।

चरघातांकी रूप में $(1+\iu)^n = 2^{n/2}\, \eu^{\iu n\pi/4}$, जो वास्तविक है
यदि और केवल यदि $\sin\frac{n\pi}{4} = 0$, अर्थात् $4 \mid n$। अतः $(1+\iu)^n
\in \R$ ठीक $4$ के गुणजों के लिए (मान $(-4)^{n/4}$)।
\end{solution}

\begin{solution}{exo:b1:complex:3}
$\abs{z - 2} = \abs{z + \iu}$: बिंदुओं $2$ और $-\iu$ से समदूरी --- अर्थात्
$(2, 0)$ और $(0, -1)$ को जोड़ने वाले रेखाखंड का लंब समद्विभाजक।

$\abs{z - 1} = 2$: केंद्र $1$ और त्रिज्या $2$ वाला वृत्त।

$\frac{z-1}{z+1} \in \iu\R$: \cref{prop:b1:complex:geometry}~(3) से, बिंदु
$M(z)$, $A(1)$, $B(-1)$ यह संतुष्ट करते हैं: रेखाएँ $MA$ और $MB$ लंब हैं
(अथवा $z = 1$, जहाँ भागफल $0 \in \iu\R$ है)। \omterm{def:b1:logic:sets}{समुच्चय} व्यास $[-1, 1]$ वाला
वृत्त (अर्थात् इकाई वृत्त) है, जिसमें से वह बिंदु $-1$ हटा दिया गया है जहाँ
भागफल अपरिभाषित है। \emph{संगणना से जाँच:} $z = \eu^{\iu\theta}$ से
$\frac{z-1}{z+1} = \frac{\eu^{\iu\theta/2}(\eu^{\iu\theta/2} -
\eu^{-\iu\theta/2})}{\eu^{\iu\theta/2}(\eu^{\iu\theta/2} +
\eu^{-\iu\theta/2})} = \iu\tan\frac\theta2 \in \iu\R$ मिलता है।
\end{solution}

\begin{solution}{exo:b1:complex:4}
रैखिकीकरण:
\[
\sin^4\theta
= \Bigl(\frac{\eu^{\iu\theta} - \eu^{-\iu\theta}}{2\iu}\Bigr)^{\!4}
= \frac{\eu^{4\iu\theta} - 4\eu^{2\iu\theta} + 6 - 4\eu^{-2\iu\theta}
+ \eu^{-4\iu\theta}}{16}
= \frac{\cos 4\theta - 4\cos 2\theta + 3}{8}.
\]
प्रसार: दे मॉयवर से $\cos 4\theta = \Re\bigl((c + \iu s)^4\bigr) = c^4 -
6c^2 s^2 + s^4$, जहाँ $c = \cos\theta$, $s = \sin\theta$; $s^2 = 1 - c^2$
प्रतिस्थापित करने पर:
\[
\cos 4\theta = c^4 - 6c^2(1 - c^2) + (1 - c^2)^2
= 8c^4 - 8c^2 + 1 .
\]
\end{solution}

\begin{solution}{exo:b1:complex:5}
$8\iu = 8\,\eu^{\iu\pi/2}$, अतः घनमूल $2\,\eu^{\iu(\pi/6 + 2k\pi/3)}$, $k =
0, 1, 2$ हैं:
\[
z_0 = 2\eu^{\iu\pi/6} = \sqrt 3 + \iu,\quad
z_1 = 2\eu^{5\iu\pi/6} = -\sqrt 3 + \iu,\quad
z_2 = 2\eu^{3\iu\pi/2} = -2\iu :
\]
त्रिज्या $2$ वाले वृत्त पर एक समबाहु त्रिभुज।

$-4 = 4\eu^{\iu\pi}$, अतः $z^4 = -4$ के हल $\sqrt 2\, \eu^{\iu(\pi/4 +
k\pi/2)}$ हैं: $\;1 + \iu$, $-1 + \iu$, $-1 - \iu$, $1 - \iu$। \omterm{def:b1:complex:field}{संयुग्मी}
मूलों को युग्मित करने पर:
\[
X^4 + 4 = \bigl(X^2 - 2X + 2\bigr)\bigl(X^2 + 2X + 2\bigr),
\]
क्योंकि $(X - (1+\iu))(X - (1-\iu)) = X^2 - 2X + 2$, और दूसरे युग्म के लिए
भी इसी प्रकार। (जाँचने के लिए प्रसार कीजिए।)
\end{solution}

\begin{solution}{exo:b1:complex:6}
$C_n + \iu S_n = \sum_{k=0}^{n} \eu^{\iu k\theta} =
\frac{\eu^{\iu(n+1)\theta} - 1}{\eu^{\iu\theta} - 1}$ (गुणोत्तर योग, अनुपात
$\eu^{\iu\theta} \neq 1$)। अंश और हर पर अर्ध-कोण गुणनखंडन
(\cref{met:b1:complex:trig}~(4)):
\[
\frac{2\iu\sin\frac{(n+1)\theta}{2}\, \eu^{\iu(n+1)\theta/2}}
{2\iu\sin\frac{\theta}{2}\, \eu^{\iu\theta/2}}
= \frac{\sin\frac{(n+1)\theta}{2}}{\sin\frac{\theta}{2}}\,
\eu^{\iu n\theta/2}.
\]
वास्तविक और काल्पनिक भाग लेने पर:
\[
C_n = \frac{\sin\frac{(n+1)\theta}{2}}{\sin\frac{\theta}{2}}
\cos\frac{n\theta}{2},
\qquad
S_n = \frac{\sin\frac{(n+1)\theta}{2}}{\sin\frac{\theta}{2}}
\sin\frac{n\theta}{2}.
\]
\end{solution}

\begin{solution}{exo:b1:complex:7}
विविक्तकर: $\Delta = (3 + 4\iu)^2 - 4(-1 + 5\iu) = 9 + 24\iu - 16 + 4 -
20\iu = -3 + 4\iu$। $-3 + 4\iu$ के वर्गमूल: $x^2 - y^2 = -3$, $2xy = 4$,
$x^2 + y^2 = 5$ हल कीजिए; फिर $x^2 = 1$, $y = 2/x$, वही चिह्न: $\delta =
\pm(1 + 2\iu)$। अतः
\[
z = \frac{(3 + 4\iu) \pm (1 + 2\iu)}{2}
\in \{\, 2 + 3\iu,\; 1 + \iu \,\}.
\]
\emph{जाँच:} योग $= 3 + 4\iu$ और गुणनफल $(2+3\iu)(1+\iu) = -1 + 5\iu$, जैसा
गुणांक चाहते हैं।
\end{solution}

\begin{solution}{exo:b1:complex:8}
\begin{enumerate}
  \item $n = 5$ के साथ \cref{prop:b1:complex:sumroots}।
  \item (1) से $u + v = \omega + \omega^2 + \omega^3 + \omega^4 = -1$।
  गुणनफल के लिए प्रसार कीजिए और घातांकों को $5$ के सापेक्ष घटाइए:
        \[
        uv = (\omega + \omega^4)(\omega^2 + \omega^3)
        = \omega^3 + \omega^4 + \omega^6 + \omega^7
        = \omega^3 + \omega^4 + \omega + \omega^2 = -1 .
        \]
        अतः $u$ और $v$ का योग $-1$ और गुणनफल $-1$ है: वे $X^2 + X - 1$ के
        दोनों मूल हैं।
  \item $u = \omega + \conj\omega = 2\cos\frac{2\pi}{5} > 0$ (कोण न्यून है),
  और $X^2 + X - 1$ का धनात्मक मूल $\frac{-1 + \sqrt 5}{2}$ है। अतः
  $\cos\frac{2\pi}{5} = \frac{\sqrt 5 - 1}{4}$।
\end{enumerate}
\end{solution}

\begin{solution}{exo:b1:complex:9}
\cref{prop:b1:complex:rules}~(4) की उपपत्ति की भाँति दोनों वर्गों का प्रसार
कीजिए:
\[
\abs{z + w}^2 = \abs z^2 + \abs w^2 + 2\Re(z\conj w),
\qquad
\abs{z - w}^2 = \abs z^2 + \abs w^2 - 2\Re(z\conj w),
\]
और जोड़ दीजिए। ज्यामितीय रूप से, $\abs{z+w}$ और $\abs{z-w}$ शीर्षों $0, z,
z+w, w$ वाले समांतर चतुर्भुज के दोनों विकर्णों की लंबाइयाँ हैं, जबकि $\abs
z$ और $\abs w$ भुजाओं की लंबाइयाँ हैं: विकर्णों के वर्गों का योग चारों
भुजाओं के वर्गों के योग के बराबर है।
\end{solution}

\begin{solution}{exo:b1:complex:10}
मान लीजिए $j = \eu^{2\iu\pi/3}$, अतः $j^2 + j + 1 = 0$ और $\eu^{\iu\pi/3} =
-j^2$, $\eu^{-\iu\pi/3} = -j$। त्रिभुज प्रत्यक्ष समबाहु है यदि और केवल यदि
$c - a = -j^2 (b - a)$, और अप्रत्यक्ष समबाहु यदि और केवल यदि $c - a = -j(b -
a)$ (\cref{ex:b1:complex:equilateral})।

$P = a + jb + j^2 c$ और $Q = a + j^2 b + jc$ पर विचार कीजिए। $1 + j^2 = -j$
और $j^3 = 1$ का प्रयोग करते हुए:
\[
c - a + j^2(b - a) = c + j^2 b - (1 + j^2)\,a = c + j^2 b + ja
= j\,(a + jb + j^2 c) = jP,
\]
अतः प्रत्यक्ष स्थिति $jP = 0 \iff P = 0$ कहती है; $j^2$ के स्थान पर $j$ रखकर
वही संगणना $c - a + j(b - a) = j^2 Q$ देती है, अतः अप्रत्यक्ष स्थिति $Q = 0$
कहती है। अतः: समबाहु (किसी भी दिशा में) $\iff PQ = 0$। $j + j^2 = -1$ के साथ
प्रसार करने पर:
\[
PQ = a^2 + b^2 + c^2 + (j + j^2)(ab + bc + ca)
= a^2 + b^2 + c^2 - (ab + bc + ca).
\]
अतः त्रिभुज समबाहु है यदि और केवल यदि $a^2 + b^2 + c^2 = ab + bc + ca$।
\end{solution}

\begin{solution}{exo:b1:complex:11}
चूँकि $\omega^k$, $k = 0, \dots, n-1$, ठीक $X^n - 1$ के $n$ मूल हैं
(\cref{thm:b1:complex:roots}), और बहुपद \omterm{def:b1:logic:inj}{एकैकी} अग्रणी गुणांक वाला है:
\[
X^n - 1 = \prod_{k=0}^{n-1} (X - \omega^k)
= (X - 1) \prod_{k=1}^{n-1} (X - \omega^k).
\]
$X - 1$ से भाग देने पर: $\;1 + X + \dots + X^{n-1} = \prod_{k=1}^{n-1} (X -
\omega^k)$। $X = 1$ पर मान लेने से $P = n$ मिलता है।

अब $1 - \omega^k = -\eu^{\iu k\pi/n}\bigl(\eu^{\iu k\pi/n} - \eu^{-\iu
k\pi/n}\bigr) = -2\iu\,\eu^{\iu k\pi/n} \sin\frac{k\pi}{n}$, अतः $P = n$ में
\omterm{def:b1:complex:field}{मापांक} लेने पर (प्रत्येक $1 \leq k \leq n-1$ के लिए $\sin\frac{k\pi}n > 0$):
\[
n = \abs P = \prod_{k=1}^{n-1} 2\sin\frac{k\pi}{n}
= 2^{n-1} \prod_{k=1}^{n-1} \sin\frac{k\pi}{n},
\qquad\text{अतः}\qquad
\prod_{k=1}^{n-1} \sin\frac{k\pi}{n} = \frac{n}{2^{n-1}} .
\]
\end{solution}

\begin{solution}{exo:b1:complex:12}
$z = 1$ हल नहीं है ($2^n \neq 0$), अतः समीकरण $\bigl(\frac{z+1}{z-1}\bigr)^n
= 1$ के तुल्य है, अर्थात् $\omega = \eu^{2\iu\pi/n}$ और $k \in \{0, \dots,
n-1\}$ के साथ $\frac{z+1}{z-1} = \omega^k$। मान $k = 0$ वर्जित है ($z + 1 =
z - 1$ असंभव है)। $1 \leq k \leq n - 1$ के लिए $z + 1 = \omega^k(z - 1)$ हल
करने पर $z(1 - \omega^k) = -1 - \omega^k$ मिलता है, अतः $\varphi =
\frac{2k\pi}{n}$ और अर्ध-कोण गुणनखंडनों $1 + \eu^{\iu\varphi} =
2\cos\frac\varphi2\, \eu^{\iu\varphi/2}$, $\eu^{\iu\varphi} - 1 =
2\iu\sin\frac\varphi2 \,\eu^{\iu\varphi/2}$ के साथ:
\[
z_k = \frac{1 + \omega^k}{\omega^k - 1}
= \frac{2\cos\frac{k\pi}{n}}{2\iu\,\sin\frac{k\pi}{n}}
= -\iu\,\frac{\cos\frac{k\pi}{n}}{\sin\frac{k\pi}{n}} ,
\]
जो दावे के अनुसार शुद्ध काल्पनिक है। \omterm{def:b1:logic:map}{प्रतिचित्रण} $t \mapsto \cos t/\sin t$
$\intoo0\pi$ पर \omterm{def:b1:logic:inj}{एकैकी} है (वह कठोरतः ह्रासमान है), अतः $n - 1$ मान $z_k$
जोड़ों में भिन्न हैं: प्रसार करने पर घात $n - 1$ वाले समीकरण ($z^n$ पद कट
जाते हैं) के ठीक यही $n - 1$ हल हैं।
\end{solution}

\begin{solution}{pb:b1:complex:1}
\textbf{1.} $j = \cos\frac{2\pi}3 + \iu\sin\frac{2\pi}3 = -\frac12 +
\iu\frac{\sqrt3}2$; $j^2 = \eu^{4\iu\pi/3} = -\frac12 - \iu\frac{\sqrt3}2 =
\conj j$; $j^3 = 1$; $1 + j + j^2 = 0$ (इकाई के घनमूलों का योग,
\cref{prop:b1:complex:sumroots}); $j^{-1} = j^2$ (क्योंकि $j \cdot j^2 =
1$)। इकाई वृत्त पर $1$, $j$, $j^2$ प्रत्यक्ष समबाहु त्रिभुज के शीर्ष हैं।

\textbf{2.} केंद्र $a$ और कोण $\theta$ का घूर्णन $a$ को अचल रखता है और $a$
से निकलने वाले प्रत्येक सदिश को $\theta$ से घुमा देता है: $r(z) - a =
\eu^{\iu\theta}(z - a)$, अर्थात् $r(z) = a + \eu^{\iu\theta}(z - a)$। $r(z)
= a + \eu^{\iu\theta}(z-a)$ और $r'(z) = a' + \eu^{\iu\theta'}(z-a')$ का
संयोजन करने पर:
\[
r' \circ r\,(z) = \eu^{\iu(\theta + \theta')} z +
\text{अचर},
\]
जो $Az + B$ रूप का \omterm{def:b1:logic:map}{प्रतिचित्रण} है, जहाँ $A = \eu^{\iu(\theta+\theta')}$ का
\omterm{def:b1:complex:field}{मापांक} $1$ है। \cref{prop:b1:complex:similitude} से वह $A \neq 1$ होने पर
कोण $\arg A = \theta + \theta'$ का घूर्णन है, और $A = 1$ होने पर, अर्थात्
$\theta + \theta' \in 2\pi\Z$ होने पर, स्थानांतरण।

\textbf{3.} $(b, c, p)$ समबाहु है यदि और केवल यदि $\abs{p - b} = \abs{c - b}
= \abs{p - c}$। $q = \frac{p - b}{c - b}$ लिखने पर पहली समता $\abs q = 1$
कहती है और दूसरी $\abs{q - 1} = 1$; दोनों मिलकर $q = \eu^{\pm\iu\pi/3}$
(वृत्तों $\abs q = 1$ और $\abs{q - 1} = 1$ के दोनों प्रतिच्छेद बिंदु
$\frac12 \pm \iu\frac{\sqrt3}2$ हैं)। अतः $p = b + \eu^{\pm\iu\pi/3}(c -
b)$। $a = 0$, $b = 1$, $c = \iu$ (प्रत्यक्ष त्रिभुज) के लिए चयन
$\eu^{-\iu\pi/3}$ देता है
\[
p = 1 + \Bigl(\tfrac12 - \iu\tfrac{\sqrt3}2\Bigr)(\iu - 1)
= \tfrac{1 + \sqrt3}2\,(1 + \iu) \approx 1.37 + 1.37\,\iu ,
\]
जो रेखा $BC$ ($x + y = 1$) के $a = 0$ से परली ओर स्थित है: बाहर की ओर। चयन
$\eu^{+\iu\pi/3}$ $p \approx -0.37 - 0.37\,\iu$ देता है, जो $a$ की ही ओर है:
भीतर की ओर।

\textbf{4.} $P = a + jb + j^2c$ को $j^2$ से गुणा कीजिए: $j^2 P = j^2 a + j^3
b + j^4 c = b + jc + j^2 a$ ($j^3 = 1$, $j^4 = j$ का प्रयोग करते हुए); और
$j$ से: $jP = ja + j^2 b + c$। ये दोनों सर्वसमिकाएँ हैं। यदि $P = 0$, तो
$j^2 P = jP = 0$: कसौटी $(b, c, a)$ और $(c, a, b)$ के लिए भी सत्य है ---
अर्थात् चक्रीय अपरिवर्त्यता (जो होनी ही चाहिए: समबाहु त्रिभुज को इससे कोई
मतलब नहीं कि पहले कौन-सा शीर्ष लिखा गया)।

\textbf{5.} यदि $a = b$: $0 = a(1 + j) + j^2 c = -j^2 a + j^2 c$ ($1 + j =
-j^2$ का प्रयोग करते हुए), अतः $c = a$। यदि $b = c$: $0 = a + b(j + j^2) = a
- b$, अतः $a = b$। यदि $a = c$: $0 = a(1 + j^2) + jb = -ja + jb$, अतः $a =
b$। हर स्थिति में तीनों बिंदु संपाती हो जाते हैं।

\textbf{6.} $(a'z + b') \circ (az + b) = a'a\,z + (a'b + b')$, जहाँ $a'a
\neq 0$: वही रूप। $z \mapsto az + b$ का प्रतिलोम $z \mapsto \frac1a z -
\frac ba$ है, जो फिर उसी रूप का है। तत्समक \omterm{def:b1:logic:map}{प्रतिचित्रण} को उदासीन अवयव मानने
पर प्रत्यक्ष समरूपताएँ संयोजन के अंतर्गत एक समूह बनाती हैं।

\textbf{7.} $f(z) = az + b$ $f(z_1) = w_1$, $f(z_2) = w_2$ को संतुष्ट करता
है यदि और केवल यदि $a z_1 + b = w_1$ और $a z_2 + b = w_2$; घटाने पर $a(z_2 -
z_1) = w_2 - w_1$, अतः
\[
a = \frac{w_2 - w_1}{z_2 - z_1} \;(\neq 0), \qquad
b = w_1 - a z_1
\]
अनिवार्य हो जाते हैं, और विलोमतः यह चयन काम करता है: अस्तित्व और अद्वितीयता।

\textbf{8.} $\abs{f(z) - f(w)} = \abs{a(z - w)} = \abs a\,\abs{z - w}$: सभी
दूरियाँ $\abs a$ से गुणित हो जाती हैं। आकृति के लिए:
\[
\frac{f(c') - f(a')}{f(b') - f(a')} = \frac{a(c' -
a')}{a(b' - a')} = \frac{c' - a'}{b' - a'} .
\]
यदि दो त्रिभुजों $(a', b', c')$ और $(a'', b'', c'')$ की आकृतियाँ बराबर हों,
तो मान लीजिए $f$ वह अद्वितीय प्रत्यक्ष समरूपता है जिसके लिए $f(a') = a''$ और
$f(b') = b''$ (प्रश्न 7); तब $(a'', b'', f(c'))$ की आकृति $(a', b', c')$ की
आकृति के बराबर है, अतः $(a'', b'', c'')$ की आकृति के भी, और आकृति पहले दो
शीर्षों से तीसरा निर्धारित कर देती है: $f(c') = c''$। विलोमतः, बराबर
आकृतियाँ ऊपर दिखाई गई अपरिवर्त्यता से निकलती हैं।

\textbf{9.} $b = f(0) = 1$; $a + b = f(1) = \iu$ से $a = \iu - 1$ मिलता है।
अनुपात $\abs a = \sqrt2$, कोण $\arg(\iu - 1) = \frac{3\pi}4$। अचल बिंदु:
\[
\zeta = \frac{b}{1 - a} = \frac1{2 - \iu} = \frac{2 + \iu}5 .
\]
अतः $f$ केंद्र $\frac{2+\iu}5$, अनुपात $\sqrt2$ और कोण $\frac{3\pi}4$ वाली
प्रत्यक्ष समरूपता है।

\textbf{10.} $\alpha + \beta = 1$ और $f(z) = az + b$ के साथ:
\[
f(\alpha a' + \beta b') = a\alpha a' + a\beta b' + b
= \alpha(a a' + b) + \beta(a b' + b)
= \alpha f(a') + \beta f(b') ,
\]
जिसकी कुंजी $b = (\alpha + \beta)b$ है। अतः मध्यबिंदु ($\alpha = \beta =
\frac12$), केंद्रक (दोहराइए) और नीचे आने वाले नेपोलियन केंद्र $\alpha b +
\beta c$ \emph{समतुल्यवर्ती} हैं: त्रिभुज का रूपांतरण उन्हें भी उसी अनुसार
रूपांतरित कर देता है। अतः समरूपता-अपरिवर्ती \omterm{def:b1:logic:statement}{कथन} को एक सुस्थापित त्रिभुज के
लिए सिद्ध कर देना सबके लिए सिद्ध कर देता है --- यही अनुमति प्रश्न 21 में
प्रयुक्त हुई।

\textbf{11.} $[b, c]$ पर बाहरी शीर्ष $p_a = b + \eu^{-\iu\pi/3}(c - b)$ है
(प्रश्न 3), अतः केंद्र है
\[
n_a = \frac{b + c + p_a}3
= \frac{2b + c + \frac{1 - \iu\sqrt3}2\,(c - b)}3
= \frac{(3 + \iu\sqrt3)\,b + (3 - \iu\sqrt3)\,c}6
= \alpha b + \beta c ,
\]
जहाँ $\alpha = \frac{3 + \iu\sqrt3}6$ और $\beta = \frac{3 - \iu\sqrt3}6 =
\conj\alpha$। भुजाओं $[c, a]$ और $[a, b]$ पर वही संगणना $n_b = \alpha c +
\beta a$ और $n_c = \alpha a + \beta b$ देती है (भूमिकाओं का चक्रीय
विस्थापन)।

\textbf{12.} $\alpha + \beta = \frac{3 + \iu\sqrt3 + 3 - \iu\sqrt3}6 = 1$।
अतः
\[
\frac{n_a + n_b + n_c}3
= \frac{(\alpha + \beta)(a + b + c)}3 = \frac{a + b + c}3 :
\]
दोनों त्रिभुजों का केंद्रक एक ही है।

\textbf{13.} $\alpha$ और $\beta$ के अनुसार समूहबद्ध कीजिए और $P = a + jb +
j^2c$ पर प्रश्न 4 की सर्वसमिकाएँ लगाइए:
\[
n_a + j n_b + j^2 n_c
= \alpha\,(b + jc + j^2 a) + \beta\,(c + ja + j^2 b)
= \alpha\,j^2 P + \beta\,j P
= (\alpha j^2 + \beta j)\,P .
\]

\textbf{14.} $j = \frac{-1 + \iu\sqrt3}2$ के साथ:
\[
\alpha j = \frac{(3 + \iu\sqrt3)(-1 + \iu\sqrt3)}{12}
= \frac{-3 + 3\iu\sqrt3 - \iu\sqrt3 - 3}{12}
= \frac{-3 + \iu\sqrt3}6 = -\beta ,
\]
अतः $\alpha j + \beta = 0$, इसलिए $\alpha j^2 + \beta j = j(\alpha j +
\beta) = 0$, और प्रश्न 13 प्रत्येक त्रिभुज $(a, b, c)$ के लिए $n_a + jn_b +
j^2n_c = 0$ दे देता है। कसौटी (प्रश्न 4--5) से केंद्र एक प्रत्यक्ष समबाहु
त्रिभुज या एक ही बिंदु बनाते हैं: यही नेपोलियन की प्रमेय है।

\textbf{15.} भीतरी शीर्ष $b + \eu^{+\iu\pi/3}(c - b)$ है, और
$\eu^{+\iu\pi/3} = \frac{1 + \iu\sqrt3}2$ के साथ प्रश्न 11 की संगणना
$\alpha$ और $\beta$ की अदला-बदली कर देती है: $m_a = \beta b + \alpha c$,
$m_b = \beta c + \alpha a$, $m_c = \beta a + \alpha b$। $Q = a + j^2b + jc$
के लिए अनुरूप सर्वसमिकाओं $b + j^2c + ja = j\,Q$ और $c + j^2a + jb = j^2 Q$
का प्रयोग करते हुए:
\[
m_a + j^2 m_b + j m_c
= \beta(b + j^2 c + ja) + \alpha(c + j^2 a + jb)
= (\beta j + \alpha j^2)\, Q = j(\beta + \alpha j)\,Q = 0 ,
\]
क्योंकि $\alpha j = -\beta$ (प्रश्न 14)। अतः भीतरी केंद्र \emph{अप्रत्यक्ष}
समबाहु कसौटी को संतुष्ट करते हैं: समबाहु, पर विपरीत दिशा वाला (अथवा एक
बिंदु)।

\textbf{16.} त्रिकी $(n_a, n_b, n_c)$ (जो प्रत्यक्ष कसौटी संतुष्ट करती है)
पर प्रश्न 5 लगाने से, दो केंद्र तभी संपाती होते हैं जब तीनों हों; और तीनों
तभी संपाती होते हैं जब \emph{दोनों} कसौटियाँ सत्य हों, अर्थात् जब अतिरिक्त
रूप से $n_a + j^2 n_b + j n_c = 0$ भी हो। सर्वसमिकाओं $b + j^2c + ja = jQ$,
$c + j^2a + jb = j^2Q$ के साथ प्रश्न 13 की भाँति संगणना कीजिए:
\[
n_a + j^2 n_b + j n_c
= \alpha\,jQ + \beta\,j^2 Q = j(\alpha + \beta j)\,Q .
\]
सीधे: $\beta j = \frac{(3 - \iu\sqrt3)(-1 + \iu\sqrt3)}{12} = \frac{-3 +
3\iu\sqrt3 + \iu\sqrt3 + 3}{12} = \frac{\iu\sqrt3}3$, अतः $\alpha + \beta j
= \frac{3 + \iu\sqrt3 + 2\iu\sqrt3}6 = \frac{1 + \iu\sqrt3}2 \neq 0$। अतः
बाहरी नेपोलियन त्रिभुज तभी अपकर्षित होता है जब $Q = a + j^2b + jc = 0$,
अर्थात् जब $(a, b, c)$ अप्रत्यक्ष समबाहु त्रिभुज हो --- उस स्थिति में सभी
``बाहरी'' रचनाएँ त्रिभुज के परिबद्ध क्षेत्र की ओर मुड़ जाती हैं और उनका
केंद्र एक ही हो जाता है।

\textbf{17.} $a = 0$, $b = 1$, $c = \iu$ के साथ:
\[
n_a = \alpha + \beta\iu = \frac{(3 + \sqrt3)(1 + \iu)}6,
\qquad
n_b = \alpha\iu = \frac{-\sqrt3 + 3\iu}6,
\qquad
n_c = \beta = \frac{3 - \iu\sqrt3}6 .
\]
वर्ग-भुजाएँ: $n_a - n_b = \frac{(3 + 2\sqrt3) + \iu\sqrt3}6$ से $\abs{n_a -
n_b}^2 = \frac{(3 + 2\sqrt3)^2 + 3}{36} = \frac{24 + 12\sqrt3}{36} = \frac{2
+ \sqrt3}3$ मिलता है; $n_b - n_c = \frac{(3 + \sqrt3)(-1 + \iu)}6$ से
$\abs{n_b - n_c}^2 = \frac{2(3 + \sqrt3)^2}{36} = \frac{24 + 12\sqrt3}{36}$;
$n_c - n_a = \frac{-\sqrt3 - \iu(3 + 2\sqrt3)}6$ से फिर वही मान। तीनों
वर्ग-भुजाएँ $\frac{2 + \sqrt3}3$ के बराबर हैं: समबाहु, जैसा वादा था।

\textbf{18.} प्रसार कीजिए:
\[
(a - b)(c - d) + (a - d)(b - c)
= (ac - ad - bc + bd) + (ab - ac - bd + cd)
= ab - ad - bc + cd ,
\]
और $(a - c)(b - d) = ab - ad - bc + cd$: बराबर।

\textbf{19.} प्रश्न 18 में \omterm{def:b1:complex:field}{मापांक} लीजिए और त्रिभुज असमिका
(\cref{prop:b1:complex:rules}~(4)) लगाइए:
\[
AC \cdot BD = \abs{(a-b)(c-d) + (a-d)(b-c)}
\leq \abs{a-b}\,\abs{c-d} + \abs{a-d}\,\abs{b-c}
= AB \cdot CD + AD \cdot BC ,
\]
जहाँ समता तभी जब दोनों योज्य $0$ से निकलने वाली एक ही किरण पर हों (त्रिभुज
असमिका की समता-स्थिति)।

\textbf{20.} अर्ध-कोण गुणनखंडन (\cref{met:b1:complex:trig}~(4)):
$\eu^{\iu\theta} - \eu^{\iu\varphi} = 2\iu\,\sin\frac{\theta -
\varphi}2\;\eu^{\iu(\theta + \varphi)/2}$, और \omterm{def:b1:complex:field}{मापांक} लेने पर एकांक \omterm{def:b1:complex:field}{मापांक}
वाले गुणनखंड समाप्त हो जाते हैं: $\abs{\eu^{\iu\theta} - \eu^{\iu\varphi}} =
2\,\abs{\sin\frac{\theta - \varphi}2}$।

\textbf{21.} $p = \eu^{\iu\theta}$, $\theta \in \intoo{2\pi/3}{4\pi/3}$ के
साथ प्रश्न 20 देता है
\[
PA = 2\,\abs{\sin\tfrac\theta2},
\quad
PB = 2\,\abs{\sin\bigl(\tfrac\theta2 - \tfrac\pi3\bigr)},
\quad
PC = 2\,\abs{\sin\bigl(\tfrac\theta2 - \tfrac{2\pi}3\bigr)} .
\]
चाप पर $\frac\theta2 \in \intoo{\pi/3}{2\pi/3}$: तब $\sin\frac\theta2 > 0$;
$\frac\theta2 - \frac\pi3 \in \intoo0{\pi/3}$, अतः दूसरी ज्या धनात्मक है;
$\frac\theta2 - \frac{2\pi}3 \in \intoo{-\pi/3}0$, अतः तीसरी ऋणात्मक है और
$PC = 2\sin\bigl(\frac{2\pi}3 - \frac\theta2\bigr)$। योग से गुणनफल:
\[
\sin\Bigl(\frac\theta2 - \frac\pi3\Bigr) +
\sin\Bigl(\frac{2\pi}3 - \frac\theta2\Bigr)
= 2\,\sin\frac\pi6\,\cos\Bigl(\frac\theta2 - \frac\pi2\Bigr)
= \sin\frac\theta2 ,
\]
अतः $PB + PC = 2\sin\frac\theta2 = PA$: यही वान स्कूटन की प्रमेय है।

\textbf{22.} $p = -1$ पर: $PA = 2$, $PB = PC = 1$, और समबाहु त्रिभुज की सभी
भुजाओं की लंबाई $\sqrt3$ है, अतः टॉलेमी के दोनों पक्ष $2\sqrt3$ कहते हैं।
बीजगणितीय रूप से, $-1 - j^2 = j$ और $1 + j = -j^2$ का प्रयोग करते हुए:
\[
(a - b)(p - c) = (1 - j)\,(-1 - j^2) = (1 - j)j = j - j^2
= \iu\sqrt3 ,
\]
\[
(a - c)(b - p) = (1 - j^2)(j + 1) = 1 + j - j^2 - j^3
= j - j^2 = \iu\sqrt3 .
\]
दोनों पद बराबर हैं, अतः उनका अनुपात $1 \in \R_{>0}$ है: यही प्रश्न 19 की
समता-स्थिति है, जो $PA \cdot BC = PB \cdot AC + PC \cdot AB$ से ठीक मेल खाती
है।

\textbf{23.} $(a, b, c) = (1, j, j^2)$ के लिए: $n_a = \alpha j + \beta j^2 =
-\beta + \beta j^2$ (प्रश्न 14) $= \beta(j^2 - 1)$; संख्यात्मक रूप से
$\beta(j^2 - 1) = \frac{(3 - \iu\sqrt3)}6 \cdot \bigl(-\frac32 -
\iu\frac{\sqrt3}2\bigr) = -1$। इसी प्रकार $n_b = \alpha j^2 + \beta = -j$ और
$n_c = \alpha + \beta j = -j^2$। $(1, j, j^2)$ का बाहरी नेपोलियन त्रिभुज
$(-1, -j, -j^2)$ है: अर्थात् मूल त्रिभुज, अपने केंद्रक $0$ से परावर्तित ---
वही आकार, आधा घुमाव घूमा हुआ। समबाहु त्रिभुज नेपोलियन रचना की एक अचल आकृति
है, कोई सिकुड़ती हुई सीमा नहीं।

\textbf{24.} (क) एकांक \omterm{def:b1:complex:field}{मापांक} वाली संख्या से गुणन का घूर्णन होना शीर्षों और
केंद्रों की रचना करता है (प्रश्न 2--3, 11) और वृत्तीय दूरियों को ज्याओं में
बदल देता है (प्रश्न 20)। (ख) समूह संरचना और आकृति-अपरिवर्त्य ने, प्रश्न 10
की समतुल्यवर्तिता के द्वारा, प्रश्न 21 में परिवृत्त को इकाई वृत्त और त्रिभुज
को $(1, j, j^2)$ तक मानकीकृत करना न्यायसंगत ठहराया। (ग) अर्ध-कोण गुणनखंडन ने
प्रश्न 20 और वान स्कूटन को पूरा करने वाले योग-से-गुणनफल पद, दोनों को शक्ति
दी।

\textbf{25.} यह शब्दकोश ज्यामितीय कथनों को ऐसी बहुपदीय सर्वसमिकाओं में बदल
देता है जिनमें हर परिकल्पना एक समीकरण है: इससे मिलता है यंत्रीकरण ---
नेपोलियन की प्रमेय $\alpha j + \beta = 0$ तक सिमट गई, $\Q(\iu\sqrt3)$ में
अंकगणित की एक ही पंक्ति; और इसकी कीमत है ज्यामितीय दृश्यता --- संगणना
प्रमाणित तो कर देती है पर यह नहीं दिखाती कि केंद्र समबाहु त्रिभुज में
\emph{क्यों} बँध जाते हैं। एक न्यायसंगत उत्तर यह है कि सर्वसमिका प्रमेय की
\emph{अनिवार्यता} समझाती है (वह $a, b, c$ में तत्समक रूप से सत्य है, अतः
विन्यास की कोई चतुराई इसमें नहीं है), जबकि चित्र उसकी \emph{सामग्री} समझाता
है। यही शब्दकोश आव्यूह के रूप में तब लौटता है जब घूर्णन
\cref{ch:b1:matrices} और \cref{ch:b1:euclid} के लांबिक $2 \times 2$ आव्यूह
बन जाते हैं, जहाँ ``$\eu^{\iu\theta}$ से गुणन'' रैखिक सममिति का आदिरूप है।
\end{solution}
